\slajdid{09_2-s040}
\begin{frame}[label=09_2-s040]
\gusttekst\setlength{\abovedisplayskip}{3pt}\setlength{\belowdisplayskip}{3pt}\setlength{\abovedisplayshortskip}{2pt}\setlength{\belowdisplayshortskip}{2pt}\begin{columns}[T,onlytextwidth]\column{.47\textwidth}
U tom slučaju potreban odnos r da bi se postavila tačka $V_S$ je
\[r=\frac{V_S-V_{Tn}}{V_{DD}-V_S+V_{Tp}}\]
Prema tome ako bi želeli da $V_S=\frac{V_{DD}}2$ onda
\[r=\frac{\frac{V_{DD}}2-V_{Tn}}{\frac{V_{DD}}2+V_{Tp}}\]
i u slučaju $V_{Tn}=-V_{Tp}$ r treba da bude 1, odnosno
\[\sqrt{\frac{k_p}{k_n}}=1\quad\Rightarrow\quad\frac{k_p}{k_n}=1\]
\[\frac{k_p}{k_n}=\frac{\mu_pC_{oxp}\frac{W_p}{L_p}}{\mu_nC_{oxn}\frac{W_n}{L_n}}\]
\column{.51\textwidth}
Podrazumeva se da su tranzistori urađeni u istoj tehnologiji $C_{oxp}=C_{oxp}$ i da je $L_p=L_n$ odnosno da se samo širinama kanala definišu odnosi (kao što smo i ranije radili) pa je potreban odnos
\[\color{DEcrvena}\frac{W_p}{W_n}=\frac{\mu_n}{\mu_p}\]
ili u opštem slučaju
\[\color{DEcrvena}\frac{W_p}{W_n}=\frac{\mu_n}{\mu_p}\left(\frac{V_S-V_{Tn}}{V_{DD}-V_S+V_{Tp}}\right)^2\]
\end{columns}
\end{frame}
